\chapter{Гладкий относительный параметр порядка и изотипическая лазейка}
\label{ch:v8-smooth-relative-background-order-parameter-gate}

\section{Расширение класса кандидатов}

Полярная часть $\operatorname{Pol}(A)$ не гладка при переходе через
$A=0$. Однако из этого ещё не следует запрет всякого гладкого относительного
параметра порядка. Физическая калибровочная группа меньше полной группы
$U(10)\times U(11)$, а потому её представление сохраняет дополнительные
изотипические проекторы.

На исходных  и целевых концах уже выведены кварк-лептонные разложения
\begin{equation}
 E_s=E_s^q\oplus E_s^\ell,
 \qquad
 E_t=E_t^q\oplus E_t^\ell,
 \qquad
 \dim(E_s^q,E_s^\ell)=(6,5),
 \quad
 \dim(E_t^q,E_t^\ell)=(6,4).
 \label{eq:v8-quark-lepton-endpoint-split}
\end{equation}
Положим
\begin{equation}
 \Gamma_s=P_s^q-P_s^\ell,
 \qquad
 \Gamma_t=P_t^q-P_t^\ell.
 \label{eq:v8-endpoint-sector-gradings}
\end{equation}
Обе градуировки коммутируют со всеми двенадцатью физическими калибровочными
генераторами. Поэтому линейное отображение
\begin{equation}
 B(A)=\Gamma_tA-A\Gamma_s
 \label{eq:v8-smooth-cross-sector-order}
\end{equation}
гладко, калибровочно эквивариантно и выделяет именно стрелки, меняющие кварк-
лептонный сектор.

\section{Положительный результат}

Полный модуль переноса имеет размерность 15 над $\mathbb C$. Его алгебра
калибровочно эквивариантных эндоморфизмов имеет размерность 13, то есть одна лишь
ковариантность не выбирает единственный параметр порядка. Тем не менее
$B$ сохраняет модуль переноса с остатком $2.8\cdot10^{-15}$, коммутирует с
калибровочным представлением с остатком $1.2\cdot10^{-15}$ и имеет комплексный ранг
6. Его спектр на модуле равен
\begin{equation}
 \{-2^{\times3},0^{\times9},2^{\times3}\}.
 \label{eq:v8-smooth-order-spectrum}
\end{equation}

Соответствующая относительная кривизна
\begin{equation}
 \mathcal R_B(A)=AA^*B(A)-B(A)A^*A
 \label{eq:v8-smooth-relative-curvature}
\end{equation}
не тождественно равна нулю. На двенадцати случайных пробах её норма лежит
между $7.15$ и $36.95$, а остаток полной калибровочной ковариантности меньше
$2.1\cdot10^{-14}$. Следовательно, полярный разрыв не является единственным
способом получить относительную кривизну.

Это существенно отличается от естественного полного унитарного
функционального исчисления. Для кандидатов вида
\begin{equation}
 B_f(A)=A f(A^*A)
 \label{eq:v8-natural-functional-order}
\end{equation}
ассоциативность даёт $AA^*B_f=B_fA^*A$, поэтому
$\mathcal R_{B_f}=0$. Ненулевая гладкая ветвь возникает именно благодаря
физическому изотипическому разложению, а не благодаря сглаживанию поляризации.

\section{Почему переход ещё не замкнут}

Поскольку $B$ линейно по $A$, кривизна $\mathcal R_B$ кубична, а действие
\begin{equation}
 S_B(A)=\|\mathcal R_B(A)\|^2
 \label{eq:v8-smooth-relative-action}
\end{equation}
однородно степени шесть:
\begin{equation}
 S_B(tA)=t^6S_B(A).
 \label{eq:v8-smooth-relative-degree-six}
\end{equation}
Поэтому в $A=0$ исчезают не только градиент, но и гессиан. Этот член не
может создать семь отрицательных запусковых мод и не заменяет квадратичный
рёберно-ходжева механизм Тома VII.

На физическом фоне $A_0$ также выполнено $B(A_0)=0$ и $S_B(A_0)=0$.
Однако линеаризация кривизны там ненулевая. Гессиан $S_B$ на 30-мерной
вещественной реаллификации модуля переноса положительно полуопределён, имеет
ранг 12, нульность 18 и положительный спектр
\begin{equation}
 12^{\times6}\cup16^{\times6}.
 \label{eq:v8-smooth-relative-vacuum-hessian-spectrum}
\end{equation}
Следовательно, новый член способен добавить жёсткость вакууму, но не
породить сам переход из нуля.

\begin{theorem}[Гладкая изотипическая относительная кривизна]
Кварк-лептонные градуировки концов задают без нового поля гладкое
калибровочно эквивариантное отображение $B(A)=\Gamma_tA-A\Gamma_s$. Полученная
кривизна $\mathcal R_B$ ненулевая на общем поле и даёт положительный
вакуумный гессиан ранга 12. Тем самым строгий запрет всех гладких замен
полярного параметра неверен.
\end{theorem}

\begin{nogocriterion}[Гладкая ветвь не является запусковым родителем]
Нельзя объявлять $S_B$ полным родительское действие: его гессиан в $A=0$ равен
нулю, на фоне $A_0$ сам параметр порядка исчезает, а 13-мерный коммутант
калибровочного представления оставляет множество других эквивариантных кандидатов.
Кварк-лептонная градуировка должна быть проверена внутри общего гессиана,
но одна ковариантность не доказывает её единственность.
\end{nogocriterion}

\begin{protocol}
\begin{enumerate}
 \item полный модуль переноса: \textbf{$15$ комплексных измерений};
 \item коммутант эквивариантных эндоморфизмов: \textbf{$13$};
 \item гладкий кандидат $B$: \textbf{получен};
 \item комплексный ранг $B$: \textbf{$6$};
 \item полная калибровочная ковариантность: \textbf{пройдена};
 \item ненулевая относительная кривизна: \textbf{получена};
 \item степень действия в нуле: \textbf{$6$};
 \item исходный гессиан: \textbf{нулевой};
 \item вакуумный гессиан: \textbf{положительный, ранг $12$};
 \item единственность кандидата: \textbf{не получена};
 \item полярный фон воспроизведён: \textbf{нет};
 \item полный родительское действие: \textbf{не получен};
 \item следующий гейт: \textbf{совместный гессиан новой кривизны}.
\end{enumerate}
\end{protocol}

Машинный сертификат сохранён в
\path{s2t_v8_smooth_relative_background_order_parameter_gate_results.json}.

\begin{thebibliography}{9}
\bibitem{v8-smooth-order-roepstorff}
G.~Roepstorff,
\emph{Superconnections and the Higgs Field},
arXiv:hep-th/9801040.
\bibitem{v8-smooth-order-baptistelli}
P.~H.~Baptistelli, M.~Manoel,
\emph{Relative Equivariants under Compact Lie Groups},
arXiv:1411.6602.
\end{thebibliography}